% An A0 poster made in OverLyX's Layout mode: every element is an object on the page.
\documentclass[final]{beamer}
\usefonttheme{professionalfonts} % keep beamer from overriding the math fonts
\usepackage[T1]{fontenc}
\usepackage[utf8]{inputenc}
\usepackage{amsmath}
\usepackage[sfdefault]{notomath} % Noto Sans for text and formulas
\usepackage[scale=0.9]{noto-mono}
\usepackage{bm}
\usepackage{graphicx}
\usepackage{qrcode}

% colours: change them here and the whole poster follows
\definecolor{main}{RGB}{2,61,107}
\definecolor{mainlight}{RGB}{0,112,180}
\definecolor{pale}{RGB}{236,242,248}
\definecolor{highlight}{RGB}{214,39,40}
\definecolor{ink}{RGB}{28,32,38}

\setbeamertemplate{navigation symbols}{}
\setbeamercolor{normal text}{fg=ink}
\setbeamercolor{itemize item}{fg=mainlight}
\setbeamertemplate{itemize item}[circle]

%% OverLyX ------------------------------------------------------------------
%% overlyx-settings: {"textclass":"beamer","papersize":"custom","paperwidth":"841mm","paperheight":"1189mm","overlyx_managed_settings":"[\"paperheight\",\"papersize\",\"paperwidth\"]"}
%% Packages and macros needed by the content (generated on every save; put your own preamble above this block).
\usepackage{tikz}
\usepackage{adjustbox}
\makeatletter
% OverLyX layout objects: positioned on the page from its top left corner (x, y, w, h;
% rotate: degrees about the centre), shown on the beamer overlay steps given by step=
\usetikzlibrary{svg.path,arrows.meta}
\pgfqkeys{/ol}{x/.store in=\ol@x,y/.store in=\ol@y,w/.store in=\ol@w,h/.store in=\ol@h,rotate/.store in=\ol@rot,
  fill/.code={\ol@setcolor\ol@fill{#1}},draw/.code={\ol@setcolor\ol@draw{#1}},color/.code={\ol@setcolor\ol@color{#1}},
  line/.store in=\ol@line,radius/.store in=\ol@radius,pad/.store in=\ol@pad,valign/.store in=\ol@valign,shape/.store in=\ol@shape,
  font/.store in=\ol@font,leading/.store in=\ol@leading,opacity/.store in=\ol@opacity,step/.store in=\ol@step,align/.store in=\ol@align,
  vb/.store in=\ol@vb,crop/.store in=\ol@crop,dash/.store in=\ol@dash,arrows/.store in=\ol@arrows,
  transition/.store in=\ol@trans,name/.store in=\ol@name,master/.store in=\ol@master,ph/.store in=\ol@ph,hide/.code={\def\ol@hide{1}},
  .unknown/.code={}}
\def\ol@reset{\def\ol@x{0mm}\def\ol@y{0mm}\def\ol@w{10mm}\def\ol@h{10mm}\def\ol@rot{0}\def\ol@fill{}\def\ol@draw{}\def\ol@color{}%
  \def\ol@line{0.4pt}\def\ol@radius{0pt}\def\ol@pad{0pt}\def\ol@valign{t}\def\ol@shape{rect}\def\ol@font{}\def\ol@leading{1.2}%
  \def\ol@opacity{1}\def\ol@step{}\def\ol@align{left}\def\ol@vb{0 0 1 1}\def\ol@crop{0 0 0 0}\def\ol@dash{}\def\ol@arrows{}\def\ol@trans{}%
  \def\ol@name{}\def\ol@master{}\def\ol@ph{}\def\ol@hide{}}
\def\ol@setcolor#1#2{\def#1{}\if\relax\detokenize{#2}\relax\else\ol@@setcolor#1#2\relax\fi}
\def\ol@@setcolor#1#2#3\relax{\ifx[#2\ol@@@setcolor#1[#3\relax\else\def#1{#2#3}\fi}
\def\ol@@@setcolor#1[#2]#3\relax{\definecolor{ol\expandafter\@gobble\string#1}{#2}{#3}\edef#1{ol\expandafter\@gobble\string#1}}
\def\ol@style{\tikzset{ol@frame/.style={}}%
  \ifx\ol@fill\empty\else\tikzset{ol@frame/.append style/.expanded={fill=\ol@fill}}\fi
  \ifx\ol@draw\empty\else\tikzset{ol@frame/.append style/.expanded={draw=\ol@draw,line width=\ol@line}}\fi
  \ifx\ol@dash\empty\else\tikzset{ol@frame/.append style/.expanded={\ol@dash}}\fi
  \ifx\ol@arrows\empty\else\tikzset{ol@frame/.append style/.expanded={\ol@arrows}}\fi}
\def\ol@only#1{\only<#1>}
% an object is drawn unless it is hidden, or a placeholder of a master page (the slides have their own boxes there)
\newif\ifol@inmaster\newif\ifol@show
\def\ol@drawn{\ol@showtrue\ifx\ol@hide\empty\else\ol@showfalse\fi\ifol@inmaster\ifx\ol@ph\empty\else\ol@showfalse\fi\fi
  \ifol@show\expandafter\@firstoftwo\else\expandafter\@secondoftwo\fi}
\def\ol@place#1{\ol@drawn{\ifx\ol@step\empty#1\else\edef\ol@tmp{\noexpand\ol@only{\ol@step}}\ol@tmp{#1}\fi}{}\ignorespaces}
\def\ol@at{\path ([xshift=\ol@x+\ol@w/2,yshift=-\ol@y-\ol@h/2]current page.north west) coordinate (ol@c);}
\newsavebox\ol@box
\def\ol@boxbegin{\pgfmathsetlengthmacro\ol@iw{\ol@w-2*(\ol@pad)}\pgfmathsetlengthmacro\ol@ih{\ol@h-2*(\ol@pad)}%
  \begin{lrbox}{\ol@box}\begin{minipage}[c][\ol@ih][\ol@valign]{\ol@iw}}
\def\ol@boxdraw{\leavevmode\begin{tikzpicture}[remember picture,overlay]\ol@at
  \begin{scope}[shift={(ol@c)},rotate=\ol@rot,opacity=\ol@opacity]
  \def\ol@tmp{ellipse}\ifx\ol@shape\ol@tmp\path[ol@frame] (0,0) ellipse [x radius=\ol@w/2,y radius=\ol@h/2];%
  \else\path[ol@frame,rounded corners=\ol@radius] (-\ol@w/2,-\ol@h/2) rectangle (\ol@w/2,\ol@h/2);\fi
  \node[anchor=center,inner sep=0pt,outer sep=0pt,transform shape] at (0,0) {\usebox\ol@box};
  \end{scope}\end{tikzpicture}}
% a text box's text is set at its natural height first, then placed in the box; that height goes to
% \jobname.olx with the class's display and list spacing (OverLyX checks its editor against the PDF)
\newwrite\ol@olx
\AtBeginDocument{\immediate\openout\ol@olx=\jobname.olx
  \begingroup\ifdefined\@listi\@listi\fi\immediate\write\ol@olx{olx params above=\the\dimexpr\abovedisplayskip\relax\space
  ashort=\the\dimexpr\abovedisplayshortskip\relax\space below=\the\dimexpr\belowdisplayskip\relax\space
  bshort=\the\dimexpr\belowdisplayshortskip\relax\space leftmargin=\the\leftmargini\space labelsep=\the\labelsep\space
  itemsep=\the\dimexpr\itemsep\relax}\endgroup}
\def\ol@tboxbegin{\pgfmathsetlengthmacro\ol@iw{\ol@w-2*(\ol@pad)}\pgfmathsetlengthmacro\ol@ih{\ol@h-2*(\ol@pad)}%
  \begin{lrbox}{\ol@box}\begin{minipage}[t]{\ol@iw}}
\def\ol@tboxend{\xdef\ol@bs{\the\dimexpr\baselineskip\relax}\par\end{minipage}\end{lrbox}%
  \ifol@inmaster\else\immediate\write\ol@olx{olx box \ifcsname c@framenumber\endcsname\the\c@framenumber\else\the\c@page\fi\space
  \ifcsname beamer@slideinframe\endcsname\the\beamer@slideinframe\else1\fi\space x=\ol@x\space y=\ol@y\space w=\ol@w\space h=\ol@h\space
  natural=\the\dimexpr\ht\ol@box+\dp\ol@box\relax\space inner=\ol@ih\space baselineskip=\ol@bs}\fi
  \sbox\ol@box{\begin{minipage}[c][\ol@ih][\ol@valign]{\ol@iw}\usebox\ol@box\end{minipage}}}
\newenvironment{olbox}[1]{\ol@reset\pgfqkeys{/ol}{#1}\ol@tboxbegin
  \ifx\ol@font\empty\else\pgfmathsetlengthmacro\ol@lead{\ol@leading*(\ol@font)}\fontsize{\ol@font}{\ol@lead}\selectfont\fi
  \ifx\ol@color\empty\else\color{\ol@color}\fi\csname ol@align@\ol@align\endcsname\ignorespaces}%
  {\ol@tboxend\ol@style\ol@place{\ol@boxdraw}}
\def\ol@align@left{\raggedright}\def\ol@align@center{\centering}\def\ol@align@right{\raggedleft}\def\ol@align@justify{}
\newenvironment{olraw}[1]{\ol@reset\pgfqkeys{/ol}{#1}\ol@boxbegin\ignorespaces}%
  {\par\end{minipage}\end{lrbox}\ol@style\ol@place{\ol@boxdraw}}
\newcommand\olshape[2]{\ol@reset\pgfqkeys{/ol}{#1}\ol@style\expandafter\ol@parsevb\ol@vb\relax
  \pgfmathsetmacro\ol@sx{(\ol@w)/\ol@vbw}\pgfmathsetmacro\ol@sy{(\ol@h)/\ol@vbh}\ol@place{\ol@shapedraw{#2}}}
\def\ol@parsevb#1 #2 #3 #4\relax{\def\ol@vbx{#1}\def\ol@vby{#2}\def\ol@vbw{#3}\def\ol@vbh{#4}}
\def\ol@shapedraw#1{\leavevmode\begin{tikzpicture}[remember picture,overlay]\ol@at
  \begin{scope}[shift={(ol@c)},rotate=\ol@rot,shift={(-\ol@w/2,\ol@h/2)},xscale=\ol@sx,yscale=-\ol@sy,shift={(-\ol@vbx pt,-\ol@vby pt)}]
  \path[ol@frame,opacity=\ol@opacity] svg {#1};\end{scope}\end{tikzpicture}}
\newcommand\olimage[2]{\ol@reset\pgfqkeys{/ol}{#1}\expandafter\ol@parsecrop\ol@crop\relax\ol@place{\ol@imagedraw{#2}}}
\def\ol@parsecrop#1 #2 #3 #4\relax{\def\ol@cl{#1}\def\ol@ct{#2}\def\ol@cr{#3}\def\ol@cb{#4}}
\def\ol@imagedraw#1{\leavevmode\begin{tikzpicture}[remember picture,overlay]\ol@at
  \node[anchor=center,inner sep=0pt,outer sep=0pt,rotate=\ol@rot,opacity=\ol@opacity] at (ol@c)
  {\adjincludegraphics[trim={\ol@cl\width} {\ol@cb\height} {\ol@cr\width} {\ol@ct\height},clip,width=\ol@w,height=\ol@h]{#1}};\end{tikzpicture}}
\newsavebox\ol@trash
\newenvironment{olgroup}[1]{\ol@reset\pgfqkeys{/ol}{#1}\ol@drawn{\ifx\ol@step\empty\def\ol@grpend{}\else
  \edef\ol@tmp{\noexpand\begin{onlyenv}<\ol@step>}\ol@tmp\def\ol@grpend{\end{onlyenv}}\fi}%
  {\begin{lrbox}{\ol@trash}\def\ol@grpend{\end{lrbox}}}\ignorespaces}{\ol@grpend\ignorespacesafterend}
% master pages: \begin{olmaster}{name=…,fill=…,master=…} keeps its objects (and keys) under its name;
% \olpage{master=…} draws them behind the frame's own — the base master's first, the placeholders not
\ExplSyntaxOn
\cs_new_protected:Npn \ol@mstore #1#2 { \tl_gclear_new:c { g__ol_master_ #1 _tl } \tl_gset:cn { g__ol_master_ #1 _tl } {#2} }
\cs_new:Npn \ol@muse #1 { \tl_if_exist:cT { g__ol_master_ #1 _tl } { \tl_use:c { g__ol_master_ #1 _tl } } }
\ExplSyntaxOff
\NewDocumentEnvironment{olmaster}{m +b}{\ol@reset\pgfqkeys{/ol}{#1}\ifx\ol@name\empty\else
  \expandafter\gdef\csname ol@mk@\ol@name\endcsname{#1}\expandafter\ol@mstore\expandafter{\ol@name}{#2}\fi}{\ignorespacesafterend}
\newcount\ol@depth
\def\ol@applykeys#1{\pgfqkeys{/ol}{#1}}
\def\ol@mkeys#1{\expandafter\let\expandafter\ol@mk\csname ol@mk@#1\endcsname\ol@reset\expandafter\ol@applykeys\expandafter{\ol@mk}}
\def\ol@masterfill#1{\ifcsname ol@mk@#1\endcsname\ol@mkeys{#1}\ifx\ol@fill\empty\ifx\ol@master\empty\else\ifnum\ol@depth<8
  \advance\ol@depth1 \edef\ol@tmp{\noexpand\ol@masterfill{\ol@master}}\ol@tmp\fi\fi\else\let\ol@pfill\ol@fill\fi\fi}
\def\ol@drawmaster#1{\ifcsname ol@mk@#1\endcsname\begingroup\ol@mkeys{#1}\ifx\ol@master\empty\else\ifnum\ol@depth<8
  \advance\ol@depth1 \edef\ol@tmp{\noexpand\ol@drawmaster{\ol@master}}\ol@tmp\fi\fi\ol@inmastertrue\ol@muse{#1}\endgroup\fi}
\newcommand\olpage[1]{\ol@reset\pgfqkeys{/ol}{#1}\let\ol@pfill\ol@fill\let\ol@pmaster\ol@master\let\ol@ptrans\ol@trans\ol@depth=0
  \ifx\ol@pfill\empty\ifx\ol@pmaster\empty\else\edef\ol@tmp{\noexpand\ol@masterfill{\ol@pmaster}}\ol@tmp\fi\fi
  \ifx\ol@pfill\empty\else\leavevmode\begin{tikzpicture}[remember picture,overlay]
  \fill[\ol@pfill] (current page.north west) rectangle (current page.south east);\end{tikzpicture}\fi
  \ifx\ol@ptrans\empty\else\ifcsname trans\ol@ptrans\endcsname\csname trans\ol@ptrans\endcsname\fi\fi
  \ifx\ol@pmaster\empty\else\ol@depth=0 \edef\ol@tmp{\noexpand\ol@drawmaster{\ol@pmaster}}\ol@tmp\fi\ignorespaces}
\@ifundefined{only}{\def\only<#1>#2{#2}\newenvironment{onlyenv}{\ol@skipspec}{}\def\ol@skipspec<#1>{}}{}
\@ifundefined{note}{\newcommand\note[2][]{}}{}
% this default might be overridden by plain title style
\newcommand\makebeamertitle{\frame{\maketitle}}%
% (ERT) argument for the TOC
\AtBeginDocument{%
  \let\origtableofcontents=\tableofcontents
  \def\tableofcontents{\@ifnextchar[{\origtableofcontents}{\gobbletableofcontents}}
  \def\gobbletableofcontents#1{\origtableofcontents}
}
\makeatother
\usepackage{geometry}
\geometry{paperwidth=841mm,paperheight=1189mm}
%% end OverLyX --------------------------------------------------------------

\begin{document}
\begin{frame}[plain]
\olshape{x=0mm,y=0mm,w=841mm,h=11mm,vb=0 0 841 11,fill=main,name=Top bar,lock}{M 0 0 L 841 0 L 841 11 L 0 11 Z}
\begin{olbox}{x=30mm,y=26mm,w=610mm,h=70mm,font=88pt,leading=1.13,color=main,grow,name=Title}
\textbf{The title of your poster says what you found, in one line}
\end{olbox}
\begin{olbox}{x=30mm,y=104mm,w=610mm,h=18mm,font=40pt,leading=1.24,color=ink,grow,name=Authors}
@@NAME@@$^{1}$\quad{}Second Author$^{2}$\quad{}Third Author$^{1}$
\end{olbox}
\begin{olbox}{x=30mm,y=126mm,w=610mm,h=34mm,font=26pt,leading=1.33,color=ink!70,grow,name=Affiliations}
\textcolor{mainlight}{$^{1}$}\,Your Institute, City, Country\qquad{}\textcolor{mainlight}{$^{2}$}\,Another Institute, City, Country

\texttt{you@example.org}
\end{olbox}
\olshape{x=668mm,y=20mm,w=1.06mm,h=136mm,vb=0 0 1.06 136,fill=main!20,name=Header divider}{M 0 0 L 1.06 0 L 1.06 136 L 0 136 Z}
\begin{olraw}{x=701mm,y=22mm,w=110mm,h=110mm,name=QR code}
{\color{main}\qrcode[height=\linewidth]{https://overlyx.app}}
\end{olraw}
\begin{olbox}{x=691mm,y=138mm,w=130mm,h=12mm,font=28pt,color=main,align=center,grow,name=QR code label}
\textbf{Paper \& code}
\end{olbox}
\olshape{x=0mm,y=168mm,w=841mm,h=3.5mm,vb=0 0 841 3.5,fill=main,name=Header rule,lock}{M 0 0 L 841 0 L 841 3.5 L 0 3.5 Z}
\olshape{x=0mm,y=171.5mm,w=841mm,h=2.5mm,vb=0 0 841 2.5,fill=mainlight,name=Header rule (light),lock}{M 0 0 L 841 0 L 841 2.5 L 0 2.5 Z}
\begin{olbox}{x=30mm,y=188mm,w=781mm,h=64.2mm,fill=pale,radius=4.94mm,name=Abstract card,lock}
\end{olbox}
\olshape{x=30mm,y=188mm,w=4.23mm,h=64.2mm,vb=0 0 4.23 64.2,fill=main,name=Abstract bar}{M 0 0 L 4.23 0 L 4.23 64.2 L 0 64.2 Z}
\begin{olbox}{x=43mm,y=195mm,w=110mm,h=16mm,font=44pt,color=main,grow,name=Abstract heading}
\textbf{Abstract}
\end{olbox}
\begin{olbox}{x=160mm,y=195mm,w=643mm,h=50.2mm,font=27pt,leading=1.32,color=ink,align=justify,grow,name=Abstract}
This poster is an example to start from, made in the \emph{Layout} mode of OverLyX. Every element on it --- this text, the headings, the coloured bands, the figures and the QR code --- is an object that you can select, move, resize and restyle, as in Keynote or Inkscape, while the text inside keeps the typography of \LaTeX , formulas included. The file is an ordinary beamer document: it compiles to an A0 PDF anywhere, and its git history stays readable. Replace the content with your own and delete what you do not need.
\end{olbox}
\begin{olbox}{x=30mm,y=306mm,w=21mm,h=21mm,fill=main,valign=c,shape=ellipse,font=44pt,color=white,align=center,name=Section 1 number}
\textbf{1}
\end{olbox}
\begin{olbox}{x=59mm,y=307.5mm,w=348.5mm,h=18mm,font=46pt,leading=1.174,color=main,grow,name=Section 1 title}
\textbf{Everything is an object}
\end{olbox}
\olshape{x=30mm,y=329.5mm,w=377.5mm,h=1.06mm,vb=0 0 377.5 1.06,fill=main!25,name=Section 1 rule}{M 0 0 L 377.5 0 L 377.5 1.06 L 0 1.06 Z}
\begin{olbox}{x=30mm,y=337mm,w=377.5mm,h=50.9mm,font=27pt,leading=1.34,color=ink,align=justify,grow,name=Section 1 text}
Click an object to select it, Shift+click to add more. Drag it to move it --- it snaps to the page, to the other objects and to their centres --- and pull its handles to resize or rotate it. A second click (or a double click) edits the text of a box; Esc returns to the box. The arrow keys nudge a selection by 1\,mm (with Shift by 10\,mm).
\end{olbox}
\begin{olbox}{x=30mm,y=393.9mm,w=377.5mm,h=40.3mm,font=27pt,leading=1.34,color=ink,grow,name=Section 1 list}
\begin{itemize}
\item \textbf{T} text box, \textbf{R} rectangle, \textbf{E} ellipse, \textbf{L} line, \textbf{A} arrow, \textbf{B} pen
\item the \emph{Layout} toolbar: fill, outline, alignment, arrangement, groups
\item pinch or Ctrl + scroll to zoom; \emph{Fit} shows the whole page
\end{itemize}
\end{olbox}
\begin{olbox}{x=30mm,y=440.2mm,w=377.5mm,h=12.8mm,font=27pt,leading=1.34,color=ink,align=justify,grow,name=Section 1 formula text}
Formulas are typed as in any OverLyX document and edited where they stand:
\end{olbox}
\begin{olbox}{x=30mm,y=459mm,w=377.5mm,h=50mm,fill=pale,radius=3.53mm,pad=2mm,valign=c,font=30pt,color=ink,name=Equation 1}
\[
\mathbf{h}_{t}=\phi\left(W\mathbf{h}_{t-1}+U\mathbf{x}_{t}\right),\qquad\hat{\mathbf{y}}_{t}=V\mathbf{h}_{t}
\]
\end{olbox}
\olshape{x=30mm,y=459mm,w=2.82mm,h=50mm,vb=0 0 2.82 50,fill=main,name=Equation 1 bar}{M 0 0 L 2.82 0 L 2.82 50 L 0 50 Z}
\begin{olbox}{x=30mm,y=575.58mm,w=21mm,h=21mm,fill=main,valign=c,shape=ellipse,font=44pt,color=white,align=center,name=Section 2 number}
\textbf{2}
\end{olbox}
\begin{olbox}{x=59mm,y=577.08mm,w=348.5mm,h=18mm,font=46pt,leading=1.174,color=main,grow,name=Section 2 title}
\textbf{Figures}
\end{olbox}
\olshape{x=30mm,y=599.08mm,w=377.5mm,h=1.06mm,vb=0 0 377.5 1.06,fill=main!25,name=Section 2 rule}{M 0 0 L 377.5 0 L 377.5 1.06 L 0 1.06 Z}
\olimage{x=58.75mm,y=614.58mm,w=320mm,h=114.35mm,name=Figure 1}{figures/model}
\begin{olbox}{x=30mm,y=736.93mm,w=377.5mm,h=21.1mm,font=23pt,leading=1.304,color=ink!70,align=justify,grow,name=Figure 1 caption}
\textbf{Figure 1.} Add images with the image tool, or drop a PDF, PNG or JPEG onto the page; double-click an image to crop it. This one is a PDF drawn with TikZ.
\end{olbox}
\begin{olbox}{x=30mm,y=824.62mm,w=21mm,h=21mm,fill=main,valign=c,shape=ellipse,font=44pt,color=white,align=center,name=Section 3 number}
\textbf{3}
\end{olbox}
\begin{olbox}{x=59mm,y=826.12mm,w=348.5mm,h=18mm,font=46pt,leading=1.174,color=main,grow,name=Section 3 title}
\textbf{Text, colours and styles}
\end{olbox}
\olshape{x=30mm,y=848.12mm,w=377.5mm,h=1.06mm,vb=0 0 377.5 1.06,fill=main!25,name=Section 3 rule}{M 0 0 L 377.5 0 L 377.5 1.06 L 0 1.06 Z}
\begin{olbox}{x=30mm,y=855.62mm,w=377.5mm,h=50.9mm,font=27pt,leading=1.34,color=ink,align=justify,grow,name=Section 3 text}
Text boxes have a base font size and line spacing, a margin, a background and an outline, and set their text ragged right, justified or centred. The colours of this poster are defined once at the top of the file (\texttt{\textbackslash definecolor}); change one there and every object that uses it follows. Named sizes such as \emph{small} and \emph{large}, bold, italics and \textcolor{highlight}{colour} work inside boxes as in documents.
\end{olbox}
\begin{olbox}{x=30mm,y=912.52mm,w=377.5mm,h=38.2mm,font=27pt,leading=1.34,color=ink,align=justify,grow,name=Section 3 more}
Objects that belong together can be grouped (Ctrl+G), and background shapes locked so that clicks go through them. The file order is the drawing order: \emph{Bring to front} and \emph{Send to back} change it.
\end{olbox}
\begin{olbox}{x=30mm,y=1017.3mm,w=21mm,h=21mm,fill=main,valign=c,shape=ellipse,font=44pt,color=white,align=center,name=Section 4 number}
\textbf{4}
\end{olbox}
\begin{olbox}{x=59mm,y=1018.8mm,w=348.5mm,h=18mm,font=46pt,leading=1.174,color=main,grow,name=Section 4 title}
\textbf{Pages, steps and notes}
\end{olbox}
\olshape{x=30mm,y=1040.8mm,w=377.5mm,h=1.06mm,vb=0 0 377.5 1.06,fill=main!25,name=Section 4 rule}{M 0 0 L 377.5 0 L 377.5 1.06 L 0 1.06 Z}
\begin{olbox}{x=30mm,y=1048.3mm,w=377.5mm,h=38.2mm,font=27pt,leading=1.34,color=ink,align=justify,grow,name=Section 4 text}
A poster is a single page; add more and the document becomes a slide deck. Every object can appear on a step of its own (\emph{Animation} on the Layout toolbar), a page can have a background colour and a transition, and speaker notes go below the page.
\end{olbox}
\begin{olbox}{x=433.5mm,y=306mm,w=21mm,h=21mm,fill=main,valign=c,shape=ellipse,font=44pt,color=white,align=center,name=Section 5 number}
\textbf{5}
\end{olbox}
\begin{olbox}{x=462.5mm,y=307.5mm,w=348.5mm,h=18mm,font=46pt,leading=1.174,color=main,grow,name=Section 5 title}
\textbf{Method}
\end{olbox}
\olshape{x=433.5mm,y=329.5mm,w=377.5mm,h=1.06mm,vb=0 0 377.5 1.06,fill=main!25,name=Section 5 rule}{M 0 0 L 377.5 0 L 377.5 1.06 L 0 1.06 Z}
\begin{olbox}{x=433.5mm,y=337mm,w=377.5mm,h=25.5mm,font=27pt,leading=1.34,color=ink,align=justify,grow,name=Section 5 text}
Say in two or three sentences what you did. A formula in its own card draws the eye to the model or the objective:
\end{olbox}
\begin{olbox}{x=433.5mm,y=368.5mm,w=377.5mm,h=50mm,fill=pale,radius=3.53mm,pad=2mm,valign=c,font=30pt,color=ink,name=Equation 2}
\[
\mathcal{L}(\theta)=\frac{1}{N}\sum_{i=1}^{N}\ell\left(f_{\theta}(\mathbf{x}_{i}),\,y_{i}\right)+\lambda\lVert\theta\rVert^{2}
\]
\end{olbox}
\olshape{x=433.5mm,y=368.5mm,w=2.82mm,h=50mm,vb=0 0 2.82 50,fill=main,name=Equation 2 bar}{M 0 0 L 2.82 0 L 2.82 50 L 0 50 Z}
\begin{olbox}{x=433.5mm,y=424.5mm,w=377.5mm,h=25.5mm,font=27pt,leading=1.34,color=ink,align=justify,grow,name=Section 5 more}
Numbered equations, cross-references and citations work as in a paper when you need them, but a poster rarely does.
\end{olbox}
\begin{olbox}{x=433.5mm,y=490.33mm,w=21mm,h=21mm,fill=main,valign=c,shape=ellipse,font=44pt,color=white,align=center,name=Section 6 number}
\textbf{6}
\end{olbox}
\begin{olbox}{x=462.5mm,y=491.83mm,w=348.5mm,h=18mm,font=46pt,leading=1.174,color=main,grow,name=Section 6 title}
\textbf{Results}
\end{olbox}
\olshape{x=433.5mm,y=513.83mm,w=377.5mm,h=1.06mm,vb=0 0 377.5 1.06,fill=main!25,name=Section 6 rule}{M 0 0 L 377.5 0 L 377.5 1.06 L 0 1.06 Z}
\begin{olbox}{x=433.5mm,y=521.33mm,w=377.5mm,h=25.5mm,font=27pt,leading=1.34,color=ink,align=justify,grow,name=Section 6 text}
Replace this plot with your own: select it and choose another image, or drop a new file onto the page and delete this one.
\end{olbox}
\olimage{x=433.5mm,y=552.83mm,w=377.5mm,h=215.41mm,name=Figure 2}{figures/results}
\begin{olbox}{x=433.5mm,y=776.24mm,w=377.5mm,h=10.6mm,font=23pt,leading=1.304,color=ink!70,align=justify,grow,name=Figure 2 caption}
\textbf{Figure 2.} Accuracy during training. The numbers are illustrative; the plot is a PDF made with pgfplots.
\end{olbox}
\begin{olbox}{x=433.5mm,y=827.17mm,w=21mm,h=21mm,fill=main,valign=c,shape=ellipse,font=44pt,color=white,align=center,name=Section 7 number}
\textbf{7}
\end{olbox}
\begin{olbox}{x=462.5mm,y=828.67mm,w=348.5mm,h=18mm,font=46pt,leading=1.174,color=main,grow,name=Section 7 title}
\textbf{Key finding}
\end{olbox}
\olshape{x=433.5mm,y=850.67mm,w=377.5mm,h=1.06mm,vb=0 0 377.5 1.06,fill=main!25,name=Section 7 rule}{M 0 0 L 377.5 0 L 377.5 1.06 L 0 1.06 Z}
\begin{olbox}{x=433.5mm,y=860.17mm,w=377.5mm,h=78mm,fill=highlight!5!white,draw=highlight,line=2.5pt,radius=3.53mm,name=Key finding card,lock}
\end{olbox}
\begin{olbox}{x=442.5mm,y=866.17mm,w=359.5mm,h=14mm,font=34pt,color=highlight,grow,name=Key finding heading}
\textbf{The one thing to remember}
\end{olbox}
\begin{olbox}{x=442.5mm,y=885.17mm,w=205mm,h=25.5mm,font=27pt,leading=1.34,color=ink,grow,name=Key finding text}
Say it in a sentence, and give the formula that carries it --- people read this card first.
\end{olbox}
\begin{olbox}{x=655.5mm,y=885.17mm,w=146.5mm,h=47mm,valign=c,font=30pt,color=ink,name=Key finding formula}
\[
\mathrm{error}\propto\frac{1}{\sqrt{N}}
\]
\end{olbox}
\begin{olbox}{x=433.5mm,y=978.5mm,w=21mm,h=21mm,fill=main,valign=c,shape=ellipse,font=44pt,color=white,align=center,name=Section 8 number}
\textbf{8}
\end{olbox}
\begin{olbox}{x=462.5mm,y=980mm,w=348.5mm,h=18mm,font=46pt,leading=1.174,color=main,grow,name=Section 8 title}
\textbf{Presenting and printing}
\end{olbox}
\olshape{x=433.5mm,y=1002mm,w=377.5mm,h=1.06mm,vb=0 0 377.5 1.06,fill=main!25,name=Section 8 rule}{M 0 0 L 377.5 0 L 377.5 1.06 L 0 1.06 Z}
\begin{olbox}{x=433.5mm,y=1009.5mm,w=377.5mm,h=38.2mm,font=27pt,leading=1.34,color=ink,align=justify,grow,name=Section 8 text}
The build button compiles the PDF for the printer: the page is exactly A0 (841\,mm $\times$ 1189\,mm, \emph{Page size} on the Layout toolbar). Press F5 to show the poster full screen and zoom in while you explain. Objects can even have animation steps, as on slides.
\end{olbox}
\begin{olbox}{x=433.5mm,y=1057.7mm,w=377.5mm,h=28.8mm,font=21pt,leading=1.3,color=ink!70,grow,name=References}
\textbf{References}

{[}1{]}~L.~Lamport. \emph{\LaTeX : A Document Preparation System}. Addison-Wesley, 1994.

{[}2{]}~T.~Tantau. \emph{The Ti}k\emph{Z and PGF Packages}. Manual, CTAN.
\end{olbox}
\olshape{x=0mm,y=1100.5mm,w=841mm,h=2.5mm,vb=0 0 841 2.5,fill=main,name=Footer rule,lock}{M 0 0 L 841 0 L 841 2.5 L 0 2.5 Z}
\olshape{x=0mm,y=1103mm,w=841mm,h=86mm,vb=0 0 841 86,fill=pale,name=Footer band,lock}{M 0 0 L 841 0 L 841 86 L 0 86 Z}
\begin{olbox}{x=30mm,y=1112mm,w=117mm,h=62mm,valign=c,font=48pt,color=main,name=Take-home heading}
\textbf{Take-home}
\end{olbox}
\olshape{x=158mm,y=1115mm,w=1.06mm,h=56mm,vb=0 0 1.06 56,fill=main!30,name=Footer divider}{M 0 0 L 1.06 0 L 1.06 56 L 0 56 Z}
\begin{olbox}{x=173.5mm,y=1112mm,w=437mm,h=62mm,valign=c,font=29pt,leading=1.345,color=ink,name=Take-home}
Keep the message of the poster to three sentences here. Everything above supports them; everything else is for the conversation in front of it.
\end{olbox}
\begin{olbox}{x=623.56mm,y=1121mm,w=187.44mm,h=56mm,fill=main,radius=4.94mm,pad=6mm,valign=c,font=40pt,leading=1.143,color=white,align=center,name=Badge}
\textbf{Made with OverLyX}
\end{olbox}
\end{frame}

\end{document}
